\begin{theorem}[Actual cycle operations and their inverses]
\label{thm:peps_actual_bidirectional_cycle_operation}
\lean{TNLean.PEPS.regularActualCycleDirectionalPermutation,
TNLean.PEPS.exists_unitary_regularActualCycleBidirectionalPermutation,
TNLean.PEPS.regularDirectedTransport_treeGauge_cycleResidual,
TNLean.PEPS.regularActualCycleDirectionalPermutation_directedTransport,
TNLean.PEPS.regularTreeCycleAssignment_directedTransport,
TNLean.PEPS.regularActualTwoChordPermutation_eq_of_ne}
\leanok
\uses{thm:peps_actual_cycle_physical_permutation}
In a chosen finite region with an actual spanning tree and root, a cycle
permutation commuting with simultaneous conjugation has one original-spin
unitary for the regular virtual action and $G$-isometric site maps. It acts on
the coordinates derived from every actual input. Its adjoint implements the
inverse permutation, with the same derived tree background and the literal
exterior retained. Both statements hold on all outgoing boundary columns
and on the actual closed contraction. A controlled two-chord permutation,
or its inverse, retains every ordered bond other than its selected second chord.
This auxiliary statement uses SCP10, lines 1765--1920 and Theorem~6.16,
lines 2271--2305; the chosen tree and cycle permutation remain explicit.
\end{theorem}
\begin{proof}\leanok
Recover the input from its derived gauge and residuals, apply the uniform
common-background operation, and use its adjoint for the inverse action.
Endpoint reconstruction gives the actual directed transports. The controlled
permutation changes only the second residual.
\end{proof}

\begin{theorem}[Two-chord transport from actual vacancy]
\label{thm:peps_actual_two_chord_transport}
\lean{TNLean.PEPS.regularActualTwoChord_forwardTransport,
TNLean.PEPS.regularActualTwoChord_reverseTransport}
\leanok
\uses{thm:peps_tree_gauge_transport,thm:peps_actual_bidirectional_cycle_operation}
Let a five-edge spanning path have successive vertices $v_0,\ldots,v_5$,
with remaining chords $v_0v_5$ and $v_1v_4$ having the same ordered orientation.
Write $D_{ij}$ for the original transport from $v_i$ to $v_j$.
If the actual loop $v_1v_2v_3v_4v_1$ is vacant, the forward operation replaces
$D_{14}$ by $D_{45}^{-1}D_{05}D_{01}^{-1}$.
If the actual loop $v_0v_1v_4v_5v_0$ is vacant, the inverse operation replaces
$D_{14}$ by $D_{34}D_{23}D_{12}$.
No gauge, residual, coefficient or state equality is assumed.
These auxiliary calculations underlie SCP10, Theorem~6.16, lines 2271--2305.
\end{theorem}
\begin{proof}\leanok
Vacancy determines the residual relation after normalization by the actual
tree gauge. The controlled permutation then determines the middle residual.
Reconstructing the directed coefficient and cancelling the tree gradients
produces the stated three-transport formulas.
\end{proof}

\begin{theorem}[Fixed native operations in both axes and both directions]
\label{thm:peps_actual_route_flux_operation}
\lean{TNLean.PEPS.torusActualRouteFluxMove,
TNLean.PEPS.torusActualRouteFluxReplacement,
TNLean.PEPS.torusActualRouteFluxUpdate,
TNLean.PEPS.torusActualRouteFluxUpdate_formula,
TNLean.PEPS.torusActualRouteVacancy_iff,
TNLean.PEPS.torusActualRouteFluxMove_transport_of_vacancy,
TNLean.PEPS.torusActualRouteFluxMove_eq_update_of_vacancy,
TNLean.PEPS.IsGIsometric.exists_unitary_torusActualRouteFluxMove,
TNLean.PEPS.IsGIsometric.exists_unitary_torusActualRouteFluxUpdate}
\leanok
\uses{thm:peps_actual_two_chord_transport}
On a finite torus with both periods at least four, choose any position and
axis. The actual translated six-site tile supplies its spanning path, root
and two chords. One original-spin unitary is chosen before every actual
bond assignment, direction, boundary configuration and target-vacancy proof.
Its forward action and adjoint action implement the actual cycle permutation
and its inverse. If the target plaquette is vacant, these operations are
exactly the literal shared-bond updates
\[
\begin{array}{c|c|c}
\text{direction}&\text{shared bond}&\text{new transport}\\\hline
\text{right}&\text{right upward side}&T L B^{-1}\\
\text{up}&\text{top rightward side}&V B L^{-1}\\
\text{left}&\text{left upward side}&T^{-1} V B\\
\text{down}&\text{bottom rightward side}&V^{-1} T L
\end{array}
\]
Here $B,V,T,L$ are the original bottom rightward, right upward,
top rightward and left upward transports of the source plaquette.
Every other ordered coefficient, including each exterior and crossing bond,
is retained. Both periodic seams are included, and the ordered coefficient
is inverted precisely when the native positive step reverses the sorted edge.
The action holds on the actual open contraction and closed state.
This is a finite-torus specialization of the elementary physical movement
in SCP10, Theorem~6.16, lines 2271--2305, not the complete four-endpoint braid
\cite{gap:rmp_peps_quantum_double_g_isometry}.
\end{theorem}
\begin{proof}\leanok
Transport the fixed six-vertex path to the actual induced tile, in either axis.
The two parallel chords have the same sorting comparison at a seam.
Apply the shared two-chord calculation, identify its vacancy loop with the
actual target plaquette, and retain every unselected coefficient.
The actual directed middle transport determines its sorted coefficient.
Apply the original-spin cycle unitary and its adjoint to the derived input
coordinates; the resulting assignment is therefore the literal update.
\end{proof}

\begin{theorem}[Four numbered vacant-target physical moves]
\label{thm:peps_vacant_plaquette_physical_move}
\lean{TNLean.PEPS.torusVacantPlaquetteMoveRegion,
TNLean.PEPS.torusVacantPlaquetteMoveTarget,
TNLean.PEPS.IsGIsometric.exists_unitary_torusVacantPlaquetteMove}
\leanok
\uses{thm:peps_actual_route_flux_operation}
For the finite-torus route of width at least eight and height at least seven,
number right, up, left and down by $0,1,2,3$. At every source position and
direction there is a fixed six-spin operation implementing the literal
update above whenever the actual clockwise target holonomy is the identity.
The selected global operator is unitary and is the identity extension of
the regional operation or its adjoint. It is chosen before the actual input,
all exterior and crossing operators, every boundary configuration and the
vacancy proof. This is the elementary-operation consequence used in SCP10,
Section~6.6, lines 2340--2423; no complete braid is asserted here
\cite{gap:rmp_peps_quantum_double_g_isometry}.
\end{theorem}
\begin{proof}\leanok
For leftward and downward motion, take the tile to the left or below the
source and use the adjoint operation. Reverse the actual vacancy loop and
substitute the literal three-transport update. Unitarity is preserved by
adjoints and by identity extension outside the six sites.
\end{proof}
